\section{Introducción}

La planta \emph{Ball on Plate} es una bola que rueda sobre un plato cuadrado de acrílico de
\SI{300}{\milli\metre} $\times$ \SI{300}{\milli\metre}. El plato se apoya en una junta universal central y se inclina en
dos ejes, $\theta_x$ y $\theta_y$, con dos servos Tower Pro MG90D que lo empujan mediante
varillas. Una cámara web cenital, a \SI{450}{\milli\metre} sobre el plato, mide la posición de
la bola. La visión corre en una PC con Linux: segmenta la bola por color, convierte píxeles a
milímetros con una homografía y estima la posición y la velocidad con un filtro de Kalman. El
controlador corre en una ESP32, que recibe el estado por el puerto serie y genera el PWM de los
servos. La variable controlada es la posición $(x, y)$ del centro de la bola, en milímetros, y
la variable manipulada es la inclinación del plato, en grados.

El problema de control consiste en llevar la bola a una posición de referencia y mantenerla ahí
sin que salga del área útil del plato, $|x|, |y| \le \SI{127}{\milli\metre}$, con una inclinación
limitada por el recorrido de los servos. A diferencia de una planta térmica, esta no tiene un
equilibrio estable: sin control, cualquier inclinación lleva la bola al borde en alrededor de un
segundo. Además, el periodo de muestreo no es libre, porque lo fija la cámara a
\SI{30}{fps}, y la cadena de medición agrega un retardo que es tiempo muerto dentro del lazo.

El enunciado del proyecto pide aplicar el flujo de diseño basado en modelos con dos
estrategias de control distintas, validar ambas hasta la simulación \emph{Model-in-the-Loop} y
llevar una de ellas hasta la implementación física. No fija requisitos numéricos, así que los de
la Tabla~\ref{tab:requisitos} los fijó el equipo como una decisión de ingeniería.

\begin{table}[!t]
\caption{Requisitos de desempeño fijados por el equipo. Ensayo: bola quieta a
\SI{100}{\milli\metre} del centro en un eje y referencia en 0.}
\label{tab:requisitos}
\centering
\footnotesize
\setlength{\tabcolsep}{4pt}
\begin{tabular}{@{}lll@{}}
\toprule
\textbf{Id} & \textbf{Requisito} & \textbf{Criterio} \\
\midrule
R1 & Tiempo de establecimiento & $t_s \le \SI{3}{\second}$ (banda del \SI{5}{\percent}) \\
R2 & Sobreimpulso & $M_p \le \SI{20}{\percent}$ \\
R3 & Error estacionario & $e_{ss} \le \SI{3}{\milli\metre}$ \\
R4 & Seguridad & $|x|, |y| \le \SI{127}{\milli\metre}$ siempre \\
Diseño & Márgenes del lazo lineal & $\mathrm{PM} \ge \SI{40}{\degree}$, $\mathrm{GM} \ge \SI{6}{\deci\bel}$ \\
\bottomrule
\end{tabular}
\end{table}

Las métricas se calculan siempre sobre la posición filtrada $\xh$ en los instantes de cuadro,
que es lo que registra la PC, para que la misma definición sirva en simulación y en el
experimento: $t_s$ es el último instante con $|\xh - r|$ mayor que el \SI{5}{\percent} del
cambio, $M_p$ el máximo sobrepaso en porcentaje del cambio y $e_{ss}$ el valor absoluto de la
media del error en el último segundo.

Como en el Proyecto~1, el objetivo no es solo que el controlador funcione en simulación, sino
recorrer el flujo completo de diseño basado en modelos \cite{mw_hil}:

\begin{center}
modelado $\rightarrow$ identificación $\rightarrow$ diseño $\rightarrow$ MIL $\rightarrow$
código $\rightarrow$ SIL $\rightarrow$ PIL $\rightarrow$ implementación física.
\end{center}

El orden real del trabajo no fue estrictamente ese. Para avanzar mientras se terminaba el
prototipo, el diseño, el MIL, la generación de código, el SIL y el PIL se hicieron primero con el
modelo físico, y la identificación con el prototipo se hizo después. Por eso la
Sección~II presenta el modelo identificado, pero los controladores de la Sección~III se
diseñaron con el modelo físico; la Sección~VII analiza esa diferencia. Las dos estrategias fueron
un PID con integración condicional y un LQR con acción integral; el PID es el que continuó hasta
la implementación.
